\section{Approach and method overview}\label{sec:introduction:approach}

This thesis adopts Design Science Research (DSR)~\cite{Hevner2004DesignScienceIS,Peffers2007DSRM} as its governing methodology.
DSR is appropriate when the research contribution is a novel artefact evaluated against a practical problem: the researcher is the designer and builder, not an observer or experimenter.
Full justification of this choice, the DSR process mapping, and the evaluation criteria and thresholds are presented in Chapter~\ref{ch:methodology}.

Three primary artefacts constitute the research contribution:
\begin{enumerate}
    \item \textit{Language specification:} defines the vocabulary, structure and meaning of a quantum-implicit programming language: what programs look like, what they mean and how they relate to each other.
    \item \textit{Formal theory:} establishes that the language is internally consistent.
    Well-formed programs never get stuck; evaluation preserves program types and the compiler's translation faithfully preserves program meaning.
    \item \textit{Compiler:} demonstrates that the language is not merely theoretical by translating programs to gate-based quantum hardware for execution.
\end{enumerate}

Together, the specification, theory and compiler constitute a complete end-to-end demonstration that a quantum-implicit programming language is formally achievable.

DSR frames the research design; it does not replace formal proof methodology.
Progress and preservation proofs and semantic preservation via simulation, drawn from programming language theory, are used to evaluate the correctness of the artefact.
These are described in Chapter~\ref{ch:methodology} and developed in Chapters~\ref{ch:pub_formal_theory} and~\ref{ch:pub_compilation}.